\documentclass[conference]{IEEEtran}
\IEEEoverridecommandlockouts

\usepackage{amsmath,amssymb}
\usepackage{booktabs}
\usepackage{cite}
\usepackage{graphicx}
\usepackage{url}
\usepackage{xcolor}
\usepackage{CJKutf8}

\newcommand{\thz}{THZ1}
\newcommand{\thp}{THP1}

\title{面向共享需量管理的多线路无接触网有轨电车停站时间与充电功率联合控制}

\author{
\IEEEauthorblockN{Jun-Jie Zhao}
\IEEEauthorblockA{School of Electrical and Computer Engineering\\
The University of Sydney\\
Sydney, Australia}
}

\begin{document}
\begin{CJK*}{UTF8}{gbsn}
\maketitle

\begin{abstract}
城市轨道交通运营商正在寻求更清洁、更灵活的供能和运营方式。无接触网有轨电车利用车载储能和车站快速充电，无需全线架设接触网，既有利于可持续交通，也能减少轨道设施对城市景观的影响。然而，现有研究大多只考虑单条线路，较少研究多条无接触网有轨电车线路在车辆到站时间不一致时如何协同运行。当多辆车同时到站充电时，运营商的总用电功率会迅速升高。停站时间还会同时影响充电时长、乘客上下车和车辆后续运行。因此，多线路运营需要实时平衡总用电功率、车辆运行所需的能量和服务效率。本文首次研究多线路无接触网有轨电车的停站时间与充电功率联合控制问题，并提出一种事件驱动的深度强化学习方法。离散事件仿真模型包含两条采用不同车载储能技术的线路。每当一辆车到站时，在线决策智能体在共享总充电功率预算下选择其停站时间和充电功率。基于广州海珠有轨电车 1 号线和黄埔有轨电车 1 号线的案例表明，在多种运行时间扰动和预算条件下，该方法能够减少同时充电造成的功率峰值和预算超限，同时保证车辆有足够能量运行并维持乘客服务水平。与固定运行、遗传算法、粒子群优化、差分进化和模型预测控制相比，该方法在管理总功率需求和实时决策方面的综合表现更好。因此，它可用于协调运行条件不断变化的多条无接触网有轨电车线路。
\end{abstract}

\begin{IEEEkeywords}
无接触网有轨电车，电气化公共交通，需求管理，强化学习，离散事件仿真，近端策略优化。
\end{IEEEkeywords}

\section{引言}
无接触网轻轨和有轨电车通过车载储能和车站充电，在不连续架设接触网的情况下提供电气化公共交通。减少接触网、支柱及其配套设施，可以降低架空设施对街道景观的视觉干扰，提升城市空间的整洁性和美观性；这一优点对于人口密集的市中心和历史街区尤其重要 \cite{guerrieri2019catenary}。无接触网轻轨和有轨电车在短暂的车站停站期间为车载储能补充能量。超级电容（SC）能够承受较高的充电功率，适合在短暂的车站停站期间快速补充能量；但这种充电方式也会使车站供电系统承受持续时间短、幅值较大的负荷。当多辆车同时到站充电时，各车辆的功率需求相互叠加，可能超过馈线和变压器的规划容量或运营商的合同需量，从而引发设备过载、电压波动和高额需量费用等问题。本文将运营商设定的总电网侧充电功率规划阈值定义为\emph{总充电功率预算}，单位为 kW。

在这一共享功率约束下，车辆的充电与停站决策存在明显的权衡。较高的充电功率有助于降低车辆后续运行中的能量不足风险，但也会加剧多车之间对有限充电功率资源的竞争；延长停站时间可以增加充电时间，并为乘客上下车提供更充分的时间，却可能扰乱行车间隔，增加乘客等待时间。因此，运营商不能孤立地确定每辆车的充电功率和停站时间，而应结合各车辆的到站状态以及系统当前剩余的可分配充电功率，对多辆车的充电与停站决策进行实时协调。



与该问题相关的研究主要包括无接触网铁路能源管理、铁路时刻表编制与重新调度、机会充电（即车辆在正常运营停站期间短时补能），以及反馈与学习控制。无接触网有轨电车依靠车载储能和多次大功率充电运行，因此车辆的运行限制同时取决于储能技术和运营条件 \cite{guerrieri2019catenary}。已有研究针对电池供能的无接触网列车优化了速度曲线 \cite{ghaviha2019speed}，但尚未协调多个共用运营商功率预算的服务在停站时间和充电功率方面的决策。

% 提前声明双栏表以使其浮动到第 2 页顶部，同时保持正文中的文献编号顺序。
\nocite{wang2019multitrain,liu2020robust,yang2020biobjective,huang2021integrated,huang2023twostage,liao2020realtime,he2020charging,hu2022joint,li2023cooperative,duan2023integrated,bi2024realtime,wu2022multidepot,li2023trajectory,schulman2017ppo,holland1975adaptation,kennedy1995pso,storn1997de}

\begin{table*}[t]
\centering
\caption{与代表性相关工作的对比。破折号表示该研究并未重点考虑相应特征。SC 和 LTO 分别表示超级电容和钛酸锂电池。}
\label{tab:positioning}
\scriptsize
\setlength{\tabcolsep}{2.5pt}
\begin{tabular}{p{2.5cm}ccccccc}
\toprule
文献 & 范围 & 决策层级 & 决策变量 & 联合服务+充电 & 异构储能 & 峰值/预算 & 真实/校准数据 \\
\midrule
Ghaviha \textit{et al.} \cite{ghaviha2019speed}
& 列车
& 轨迹
& 列车运行速度
& --
& 电池
& --
& 校准 \\

Wang and Goverde \cite{wang2019multitrain}
& 多列车
& 时刻表
& 多列车运行轨迹
& --
& --
& --
& 运行 \\

Huang \textit{et al.} \cite{huang2021integrated}
& 多线路
& 时刻表
& 列车时刻表
& --
& --
& --
& 案例数据 \\

Hu \textit{et al.} \cite{hu2022joint}
& 公交网络
& 调度
& 车辆位置+充电决策
& 是
& 电池
& 电网/不确定性
& 案例数据 \\

Duan \textit{et al.} \cite{duan2023integrated}
& 公交网络
& 调度
& 服务调度+充电计划
& 是
& 电池
& 充电计划
& 案例数据 \\

Bi \textit{et al.} \cite{bi2024realtime}
& 公交线路
& 实时
& 闪充控制
& 仅联合服务
& 电池
& 充电器
& 案例数据 \\

本文研究
& 两条无接触网有轨电车
& 事件
& 停站时间+充电功率
& 是
& SC+LTO
& 共享软预算
& 案例数据 \\
\bottomrule
\end{tabular}
\end{table*}

多列车轨迹和时刻表研究会协调运行时间、行车间隔和再生能量交换 \cite{wang2019multitrain,liu2020robust,yang2020biobjective}。其他研究把时刻表优化扩展到互联线路和两阶段城市轨道交通规划 \cite{huang2021integrated,huang2023twostage}，或使用深度学习在随机扰动下实时调整列车运行 \cite{liao2020realtime}。这些工作主要控制列车运行或再生能量利用，尚未同时考虑到站时的停站与充电决策、不同的车载储能技术，以及多线路共享的软性功率预算。

与本问题较为接近的电动公交研究会联合优化充电时间、充电设施位置、车辆调度和时刻表调整 \cite{he2020charging,hu2022joint,li2023cooperative,duan2023integrated}。也有研究使用深度强化学习控制实时闪充 \cite{bi2024realtime}，或在多车场调度中考虑电网条件 \cite{wu2022multidepot}。这些研究说明服务与能源决策需要协调，但公交车场充电、充电设施选址和车辆排班，与无接触网有轨电车逐站充电仍有明显差别。

综合来看，现有研究尚未解决具有异构车载储能的多条无接触网有轨电车线路如何在共享需量预算下协调异步发生的停站与充电决策。若对各线路分别进行控制，就难以充分考虑跨线路充电时段的重叠，也无法利用其他线路暂时闲置的功率容量。表~\ref{tab:positioning} 将本文问题与代表性相关工作进行了对比。

针对上述研究空白，本文提出一种面向多线路无接触网有轨电车的事件驱动深度强化学习方法。离散事件仿真用于刻画车辆异步到站和跨线路充电重叠。每当车辆到站时，在线决策智能体根据该车状态和两条线路当前的充电情况，在共享功率需量预算下联合确定其停站时间和充电功率。案例考察两条采用不同车载储能系统、在同一个模拟早高峰内运行的无接触网有轨电车线路。案例基于广州海珠有轨电车 1 号线（\thz）和广州黄埔有轨电车 1 号线（\thp）。前者在车站为车载超级电容充电，后者同时使用超级电容和钛酸锂电池(LTO) \cite{haizhu_wiki,haizhu_baike,huangpu_wiki,huangpu_sasac}。两条线路由同一运营商统一管理，并共同受总充电功率预算约束。因此，在任一时刻，需要将两条线路上所有正在充电车辆的功率相加，并根据所得总充电功率判断是否超出预算。总充电功率预算是运营商在供电规划中设定或通过合同购买的惩罚型功率目标，短时超出会产生惩罚。预算容量始终在线路间动态共享，并根据两条线路的实时充电需求协调分配。


本文有三项贡献。第一，据我们所知，本文首次针对多条无接触网有轨电车线路，建立停站时间与充电功率一体化控制问题。该模型在运营商级总功率需量约束下，将异步车辆运行、乘客服务决策和异构车载储能系统的充电需求进行耦合，把停站与充电协同从单线路扩展到多线路场景。第二，本文建立了一个基于统一双线路离散事件仿真和马尔可夫决策过程（MDP）的事件驱动决策框架。两条线路的到站事件按统一时序处理，使每次决策能够同时考虑车辆与乘客状态、牵引能量需求、运营扰动和跨线路的实时充电活动。第三，本文针对所提出的多线路环境开发了一个基于深度强化学习（DRL）的在线决策智能体。通过面向问题设计观测空间、连续动作空间和奖励函数，该智能体能够在每次车辆到站时自适应地确定停站时间和充电功率，并协调两条线路的运输服务与能量需求。

\section{多线路无接触网有轨电车系统建模}
\subsection{双线路系统与异构储能配置}
逐次到站控制策略的性能不能仅依据静态的线路和车辆参数进行评估，因为每一次停站决策都会改变车辆离站时间、站台候车人数和车载剩余能量，并进一步影响后续车辆充电时段的重叠关系。为此，本文建立一个统一的离散事件仿真器。线路集合记为 $\mathcal{L}$；对任一线路 $\ell\in\mathcal{L}$，其车辆集合和站台集合分别记为 $\mathcal{V}_{\ell}$ 和 $\mathcal{N}_{\ell}$。各线路可以采用仅含超级电容的储能系统，也可以采用超级电容与电池组成的混合储能系统。线路、车队和储能参数均作为仿真器输入，具体案例取值见第~IV节。

在有限仿真时域 $[t_0,t_f]$ 内，各线路的车辆循环运行。车辆初始位置、车队规模、区间运行时间、在线决策智能体选择的停站时间、折返时间和运行扰动共同生成实际到站序列和行车间隔。因此，控制动作会通过后续到站时刻继续影响行车间隔、乘客积累和跨线路充电重叠。本文将 $[t_0,t_f]$ 内的一次完整仿真称为一个\emph{回合}。

\subsection{事件驱动运行过程}
在线决策智能体只在车辆到站时作出决策，但由决策产生的停站和充电时段会在相邻到站事件之间持续。因此，仿真器采用由车辆到站触发的离散事件机制。每个事件 $k$ 对应于电车 $i$ 在时间 $t_k$ 到达站台 $n$。两条线路的到站事件进入同一个全局优先队列，并按事件时间依次处理；如果多个事件同时发生，则按线路代码、车站编号和车辆编号确定处理顺序。每次到站服务完成后，仿真器更新车辆、站台、储能和运行状态，再将该车辆的下一次到站事件加入队列。因此，两条线路虽然具有不同的车辆配置和运行周期，但始终在统一的仿真时钟下演化。

车辆在车站停留期间的充电区间记为 \([t_k,t_k+d_k)\)，其中，停站时间 \(d_k\) 由第 III 节的联合决策模型确定。采用左闭右开区间意味着车辆从到站时刻 \(t_k\) 开始充电，并在时刻 \(t_k+d_k\) 结束充电；充电结束后，该车辆不再计入该时刻的总充电功率。当多辆车在同一时刻到达时，模型按照既定顺序依次处理各车辆的到站事件。每处理一辆车，其确定的充电功率都会立即计入当前总充电功率，因此，后续处理的车辆能够基于此前车辆已经形成的充电负荷作出决策。通过上述处理方式，模型可以准确表示不同线路车辆在时间上的充电重叠及其功率叠加。

\subsection{运行时间与乘客动态}
车辆停站时间和随机运行延误共同决定其后续到站时刻，并进一步影响不同车辆充电时段的重叠关系。仿真器按下式更新离站时间和下一次到站时间：
\begin{align}
    t^{\mathrm{dep}}_k &= t_k + d_k, \\
    t^{\mathrm{next}}_k &= t^{\mathrm{dep}}_k + T^{\mathrm{run}}_{k},
\end{align}
其中，事件 $k$ 表示车辆 $i$ 在时刻 $t_k$ 到达当前站台；$d_k$ 是该事件的停站时间；$t^{\mathrm{dep}}_k$ 是车辆离开当前站台的时刻；$t^{\mathrm{next}}_k$ 是事件 $k$ 生成并加入全局队列的该车辆下一站到达时刻，它不一定等于全局第 $k+1$ 个事件的发生时刻；$T^{\mathrm{run}}_k$ 是从当前站到下一站的实际区间运行时间。上述时间变量的单位均为秒。

为描述随机延误，设 $\overline{T}^{\mathrm{run}}_k$ 为计划区间运行时间，$z_k$ 为延误指示变量，则
\begin{equation}
 T^{\mathrm{run}}_k=\overline{T}^{\mathrm{run}}_k
 \left[1+(\mu_{\mathrm d}-1)z_k\right],
 \qquad z_k\sim\operatorname{Bernoulli}(p_{\mathrm d}).
\end{equation}
其中，$p_{\mathrm d}\in[0,1]$ 是单个区间发生延误的概率，$\mu_{\mathrm d}\geq1$ 是发生延误时相对于计划运行时间的放大倍数。因此，$z_k=0$ 时实际运行时间等于计划值，$z_k=1$ 时实际运行时间为 $\mu_{\mathrm d}\overline{T}^{\mathrm{run}}_k$。该式只定义一般性的运行时间转移机制；稳健性实验采用的参数组合和随机数设置在第~IV节说明。

停站时间也是车辆能够用于乘客服务的时间，因此为了减少充电重叠而缩短停站，可能使更多乘客留在站台。令 $\mathcal{N}=\bigcup_{\ell\in\mathcal{L}}\mathcal{N}_{\ell}$ 为全部站台的集合，$Q_n(t)$ 为时刻 $t$ 在站台 $n\in\mathcal{N}$ 候车的乘客数，则一个回合内的乘客总等待时间为
\begin{equation}
 W=\sum_{n\in\mathcal{N}}\int_{t_0}^{t_f}Q_n(t)\,\frac{dt}{60},
 \qquad [W]=\text{乘客-分钟}.
\end{equation}
其中，$t_0$ 和 $t_f$ 分别是回合的起止时刻；积分给出乘客等待秒数，因而除以 60 后得到乘客-分钟。在两个事件之间，$Q_n(t)$ 按给定的乘客到达率曲线增长。车辆到站后，乘客先下车；随后，车辆在剩余容量和上车人数上限 $d_k/\delta$ 内接载已经在站台候车的乘客，其中 $\delta$ 是每位乘客的假设上车服务时间，单位为秒/人，因而 $d_k/\delta$ 表示停站时间允许的最大上车人数。车辆停站期间才到达站台的乘客需要等待下一班车。上述简化规则定义了仿真中的乘客服务过程。仿真器精确计算每个事件间隔和停站时段内分段线性候车人数曲线下的面积。状态更新、等待成本、仿真结束时仍在候车的乘客数和无法上车人数均采用同一规则。

\subsection{充电与异构储能动态}
只有结合车辆离站后所需的能量，才能判断当前充电决策是否可行。因此，仿真器使用区间级能量需求模型评价每个动作。该模型为每个区间建立三角形或梯形速度曲线 \cite{ghaviha2019speed}，并使用国际单位制对加速、匀速和制动阶段进行积分。模型计入加速和匀速时的牵引损耗，并扣除再生制动回收的能量。辅助能耗取决于运行时间和车内乘客数：
\begin{equation}
    W^{\mathrm{aux}}_k = \left(P_{\mathrm{ic}} + P_{\mathrm{oc}} q_{i,k}\right)\frac{T^{\mathrm{run}}_k}{3600}.
\end{equation}
其中，$W^{\mathrm{aux}}_k$ 是事件 $k$ 后续区间的辅助能量，单位为 kWh；$P_{\mathrm{ic}}$ 是与乘客数无关的辅助功率，单位为 kW；$P_{\mathrm{oc}}$ 是每位车内乘客增加的辅助功率，单位为 kW/人；$q_{i,k}$ 是车辆 $i$ 在该区间内的车内乘客数；$T^{\mathrm{run}}_k/3600$ 将以秒计的运行时间换算为小时。车辆通过该区间所需的总能量为 $W_k = W^{\mathrm{trac}}_k + W^{\mathrm{aux}}_k$，其中 $W_k$、$W^{\mathrm{trac}}_k$ 和 $W^{\mathrm{aux}}_k$ 分别表示总能量、牵引能量和辅助能量，单位均为 kWh。

在事件 $k$，令 $P_k$ 表示分配给到站车辆的电网侧充电功率。模型假设该功率在充电区间 $[t_k,t_k+d_k)$ 内保持不变，并忽略充电连接和断开所需的时间。因此，车辆停站期间从电网获得的能量为
\begin{equation}
 E^{\mathrm{grid}}_k
 =\frac{1}{3600}\int_{t_k}^{t_k+d_k}P_k\,dt
 =\frac{P_kd_k}{3600},
\end{equation}
其中，$P_k$ 的单位为 kW，$d_k$ 的单位为秒，除以 3600 后得到的电网侧能量单位为 kWh。该电网侧能量按各储能部件的充电效率转换为实际存入的能量。对于只配备超级电容的车辆，充电量不能超过超级电容的剩余容量。对于混合储能车辆，固定的车载能量管理系统（EMS）规则会先在荷电状态限制内为超级电容充电，剩余可充能量再进入电池，同时满足电池的荷电状态和充放电倍率（C-rate）限制；C-rate 表示充放电功率相对于储能部件额定能量容量的归一化倍率。车辆牵引时也优先使用超级电容，不足部分由电池提供。如果两种储能都无法满足需求，剩余部分记为未满足牵引能量。该固定优先级规则直接确定功率在超级电容和电池之间的分配。

\section{事件驱动停站--充电联合优化方法}
\subsection{事件驱动决策问题}
在第~II节系统动态的基础上，本文把双线路协同控制表示为有限时域的事件驱动马尔可夫决策过程（MDP）。仿真器的完整状态 $s_k$ 包含全局事件队列、全部车辆与站台状态、正在进行的充电时段、运行方向和累计成本。每当任一线路的车辆到站时，系统进入一个新的决策时点；智能体作出当前车辆的停站与充电决策后，环境依次完成乘客服务、储能更新和区间运行，并推进到全局队列中的下一次到站事件。因此，决策间隔等于相邻到站事件之间经过的实际时间。

在完整状态层面，$s_k$ 和当前动作共同确定下一状态的概率分布。但策略只接收下文定义的观测 $o_k$，因此在线决策智能体面对的是部分可观测问题。控制目标是在共享功率预算下协调两条线路的异步决策，在减少功率超限和不必要充电的同时，抑制未满足牵引能量和乘客等待成本。

\subsection{观测空间与动作空间}
在线控制时，策略不能直接使用仿真器内部的全部变量，因此需要用一组紧凑变量描述到站车辆、所在站台和当前充电活动。作出动作前，智能体会获得以下 13 个观测量：
\begin{align}
 o_k = (&t_k,i,n,h_k,E^{\mathrm{sc}}_{i,k},E^{\mathrm{bat}}_{i,k},
 q_{i,k},w_{n,k},\nonumber\\
 &P^{\mathrm{tot}}_k,P^{\mathrm{own}}_k,N^{\mathrm{ch}}_k,
 E^{\mathrm{rem}}_k,\ell_k),
\end{align}
其中，$o_k$ 是事件 $k$ 的策略观测；$i$ 和 $n$ 分别是当前到站车辆和站台的编号，$\ell_k$ 是该车辆所属线路的标识。$t_k$ 是到站时刻，$h_k$ 是同一线路、同一站台和同一方向的上一辆车离站后经过的时间，单位均为秒。$E^{\mathrm{sc}}_{i,k}$ 和 $E^{\mathrm{bat}}_{i,k}$ 分别是车辆 $i$ 的超级电容和电池在事件 $k$ 的储能量，单位为 kWh；未配置电池时，$E^{\mathrm{bat}}_{i,k}=0$。$q_{i,k}$ 是车内乘客数，$w_{n,k}$ 是站台 $n$ 的候车人数。$P^{\mathrm{tot}}_k$ 是分配当前动作之前两条线路正在使用的总充电功率，$P^{\mathrm{own}}_k$ 是其中属于当前车辆所在本线路的功率，单位均为 kW。$N^{\mathrm{ch}}_k$ 是正在充电的车辆总数，$E^{\mathrm{rem}}_k$ 是这些车辆按既有充电计划尚需从电网获得的能量，单位为 kWh。车辆编号和站台编号分别在线路内部定义。

智能体在事件 $k$ 输出二维连续动作
\begin{equation}
    a_k = (d_k, \widehat{P}_k),
\end{equation}
其中，$d_k$ 是停站时间，$\widehat{P}_k$ 是智能体请求的电网侧充电功率。一般动作范围为
\begin{equation}
    d^{\min} \le d_k \le d^{\max}, \qquad
    0 \le \widehat{P}_k \le P^{\max}.
\end{equation}
其中，$d^{\min}$ 和 $d^{\max}$ 是允许的最短和最长停站时间，$P^{\max}$ 是智能体可请求的最大充电功率。仿真器根据车站功率、储能剩余容量、充电效率、荷电状态和 C-rate 限制，将 $\widehat{P}_k$ 修正为实际执行功率 $P_k\leq\widehat{P}_k$。没有充电设施的车站强制令 $P_k=0$；请求值低于不充电阈值 $P^{\mathrm{off}}$ 时也令 $P_k=0$。这些边界和阈值的案例取值见第~IV节。

对于正在进行的充电时段 $r$，设其实际充电功率为 $P_r$、开始时间为 $\tau_r$、停站时间为 $d_r$，则事件时刻 $t$ 尚需从电网获得的能量为
\begin{equation}
 E^{\mathrm{rem}}(t)=\sum_{r\in\mathcal{A}(t)}
 \frac{P_r}{3600}\left[d_r-\min\{\max(t-\tau_r,0),d_r\}\right],
\end{equation}
其中，$r$ 是充电时段的索引，$\mathcal{A}(t)$ 是时刻 $t$ 仍在进行的充电时段集合；方括号内的量是时段 $r$ 的剩余充电秒数。$P_r$ 的单位为 kW，$\tau_r$、$d_r$ 和 $t$ 的单位为秒，除以 3600 后，$E^{\mathrm{rem}}(t)$ 的单位为 kWh。该量表示现有充电计划尚需从电网获得的能量，观测中的 $E^{\mathrm{rem}}_k$ 即 $E^{\mathrm{rem}}(t_k)$。

\subsection{共享功率预算与成本函数}
预算 $B$ 是一个惩罚型规划目标，表示合同约定的需量上限或运营商购买的功率容量。系统暂时超过 $B$ 时会产生惩罚。不同的重叠充电可能以不同幅度、持续不同时间超过预算，因此单个采样功率或峰值不能完整表示超限程度。用所有充电开始和结束时刻把 $[t_0,t_f)$ 划分成连续区间，并把这些区间的集合记为 $\mathcal{I}$。每个区间 $I\in\mathcal{I}$ 内的总充电功率保持不变。因此，共享预算下的功率超限量和成本按实际时间对超额功率的平方进行积分：
\begin{align}
 O_{\mathrm{shared}}(B)&=\sum_{I\in\mathcal{I}}\frac{|I|}{3600}
 [P^{\mathrm{tot}}_I-B]_+^2,\\
 C^{\mathrm{ol}}_{\mathrm{shared}}&=c_oO_{\mathrm{shared}}(B),
\end{align}
其中，$O_{\mathrm{shared}}(B)$ 是共享预算下的功率超限量，$P^{\mathrm{tot}}_I$ 是区间 $I$ 内的总充电功率，$[x]_+=\max(0,x)$，$|I|$ 是区间持续时间，单位为秒。因此，$O_{\mathrm{shared}}$ 的单位为 $\mathrm{kW}^2\mathrm{h}$。$c_o$ 是功率超限成本权重，$C^{\mathrm{ol}}_{\mathrm{shared}}$ 是相应成本。平方项使幅度更大的超限获得更高权重，区间长度则保留其持续时间。为了比较灵活共享和固定预分配两种方式，分线路预算采用
\begin{align}
 O_{\mathrm{split}}(B)&=\sum_{\ell\in\mathcal{L}}\sum_{I\in\mathcal{I}}
 \frac{|I|}{3600}[P_{\ell,I}-B_\ell]_+^2,\\
 C^{\mathrm{ol}}_{\mathrm{split}}&=c_oO_{\mathrm{split}}(B),
\end{align}
其中，$O_{\mathrm{split}}(B)$ 是分线路预算下的功率超限量；$\ell$ 是线路索引，$P_{\ell,I}$ 是线路 $\ell$ 在区间 $I$ 内的充电功率，$B_\ell$ 是分配给该线路的预算，并满足 $\sum_{\ell\in\mathcal{L}}B_\ell=B$；$C^{\mathrm{ol}}_{\mathrm{split}}$ 是相应成本。具体的预算分配方式在第~IV节给出。因此，功率超限量是超额功率平方对实际持续时间的积分。

如果只惩罚功率超限，智能体即使在后续区间需要能量时也会倾向于不充电。因此，目标函数还必须考虑常规充电能量、未满足牵引能量和乘客等待。模型另设一个防止不必要充放电循环的累计超额项，本文称其为\emph{过充}；它是建模惩罚，不表示电池发生过电压。常规充电成本为
\begin{equation}
 C^{\mathrm{ch}}_k = c_e E^{\mathrm{ch}}_k,
\end{equation}
其中，$E^{\mathrm{ch}}_k$ 是事件 $k$ 实际存入车载储能的能量，单位为 kWh；$c_e$ 是常规充电成本权重；$C^{\mathrm{ch}}_k$ 是事件充电成本。令 $\mathcal{J}_i$ 为车辆 $i$ 的储能部件集合，$Q^{\mathrm{ch}}_{i,j,k}$ 和 $Q^{\mathrm{cons}}_{i,j,k}$ 分别为截至事件 $k$，车辆 $i$ 的部件 $j\in\mathcal{J}_i$ 累计充入和累计消耗的能量，单位均为 kWh。定义累计充电超额量
\begin{equation}
 X_{i,j,k}=\max\!\left(0,Q^{\mathrm{ch}}_{i,j,k}
 -1.25Q^{\mathrm{cons}}_{i,j,k}\right).
\end{equation}
事件 $k$ 的过充成本采用累计超额量的有符号变化值：
\begin{equation}
 C^{\mathrm{oc}}_k=c_p\sum_{j\in\mathcal{J}_i}
 \left(X_{i,j,k}-X_{i,j,k^-}\right),
\end{equation}
其中，$X_{i,j,k}$ 的单位为 kWh，$k^-$ 表示事件 $k$ 发生前一刻，$c_p$ 是过充成本权重，$C^{\mathrm{oc}}_k$ 是事件 $k$ 的过充成本。系数 1.25 是为减少不必要充放电循环而设定的固定设计裕量。车辆随后消耗能量并降低累计超额量时，负的事件成本 $C^{\mathrm{oc}}_k$ 会抵消等量的先前惩罚。在完整仿真回合内，各事件增量之和为非负量 $c_p\sum_i\sum_{j\in\mathcal{J}_i}X_{i,j,K_i}\geq0$，其中 $K_i$ 是车辆 $i$ 的最后一个事件序号。混合 EMS 在计入效率，并满足两种储能的荷电状态和充放电倍率限制后，把事件 $k$ 中仍然无法供应的牵引能量记为 $D_k\geq0$，单位为 kWh。在整个仿真回合内，
\begin{equation}
 C^{\mathrm{em}}=c_m\sum_kD_k,\qquad C^{\mathrm{wt}}=c_w W.
\end{equation}
其中，$c_m$ 和 $c_w$ 分别是未满足牵引能量和乘客等待成本权重，$C^{\mathrm{em}}$ 和 $C^{\mathrm{wt}}$ 是相应的回合总成本。事件级未满足能量成本为 $C^{\mathrm{em}}_k=c_mD_k$。

\subsection{优化目标与奖励函数}
充电、过充和未满足能量成本在相应事件发生时更新；功率超限和乘客等待成本则按实际时间积分，并将增量分配到相邻决策事件。事件 $k$ 的总成本及奖励定义为
\begin{equation}
 C_k = C^{\mathrm{ch}}_k + C^{\mathrm{oc}}_k + \Delta C^{\mathrm{ol}}_k
 + C^{\mathrm{em}}_k + \Delta C^{\mathrm{wt}}_k,
 \qquad R_k=-C_k,
\end{equation}
其中，$C_k$ 和 $R_k$ 分别是事件 $k$ 的总成本和奖励；$\Delta C^{\mathrm{ol}}_k$ 与 $\Delta C^{\mathrm{wt}}_k$ 分别是从上一决策时点到当前决策时点新增的功率超限成本和乘客等待成本。各增量项把整个仿真回合的积分分配到不同事件，避免重复计算同一时段。学习策略 $\pi$ 采用近端策略优化（PPO）实现，并最大化按事件序号折扣的期望累计奖励：
\begin{equation}
 \max_{\pi}\ J_{\gamma}(\pi)
 =\mathbb{E}_{\pi}\!\left[\sum_{k=0}^{K-1}\gamma^k R_k\right].
\end{equation}
其中，$J_{\gamma}(\pi)$ 是策略 $\pi$ 的期望折扣回报，$\mathbb{E}_{\pi}$ 对策略动作和随机运行扰动取期望，$\gamma$ 是事件级折扣因子，$K$ 是一个仿真回合内的决策事件总数，指数 $k$ 是决策事件索引。

系数 $c_e,c_o,c_m,c_p$ 和 $c_w$ 是具有相应量纲的固定设计权重，以任意目标函数单位计量。其具体取值及变化方案在第~IV节给出。

\subsection{基于 PPO 的求解方法}
本文使用近端策略优化（PPO）求解。PPO 的裁剪重要性采样比目标可以重复利用已采集的轨迹，同时降低优化过程对大幅概率比变化的激励 \cite{schulman2017ppo}。参数为 $\theta$ 的 Actor $\pi_\theta$ 在无约束潜在动作 $u_k\in\mathbb{R}^2$ 上定义因子化高斯分布：
\begin{equation}
 u_k\sim\mathcal{N}\!\left(\mu_\theta(o_k),
 \operatorname{diag}(\sigma_\theta^2)\right).
\end{equation}
其中，$\mu_\theta(o_k)\in\mathbb{R}^2$ 是随观测变化的均值向量，$\sigma_\theta\in\mathbb{R}_{>0}^2$ 是由 Actor 参数化的正标准差向量；对角协方差表示两个潜在动作分量在给定观测后相互独立。参数为 $\phi$ 的 critic $V_\phi$ 根据 $o_k$ 估计折扣累计奖励。

令 $a^{\min}=(d^{\min},0)^\top$、$a^{\max}=(d^{\max},P^{\max})^\top$、$\mathbf{1}=(1,1)^\top$。潜在动作先经过双曲正切压缩，再逐元素映射到停站时间和请求功率的盒约束允许范围：
\begin{equation}
 a_k=a^{\min}+\frac{\tanh(u_k)+\mathbf{1}}{2}
 \odot(a^{\max}-a^{\min}),
\end{equation}
其中，$\odot$ 表示逐元素乘法。该映射生成满足动作边界的策略建议值 $a_k=(d_k,\widehat{P}_k)$；仿真器随后按照第~III-B节的储能和充电设备限制，将请求功率转换为实际施加的充电功率。

训练过程中，一条 \emph{rollout} 是在一个完整仿真回合内采集的事件序列。生成 rollout 时，Actor 和 critic 参数分别固定为 $\theta_{\mathrm{old}}$ 和 $\phi_{\mathrm{old}}$；因此，$\pi_{\theta_{\mathrm{old}}}$ 是生成其中动作及所存对数概率的采样策略。

有界动作坐标中的概率密度由变量变换公式得到。对于动作维度 $j\in\{1,2\}$，
\begin{align}
 \log\pi_\theta(a_k\mid o_k)
 =&\ \log\mathcal{N}\!\left(u_k;\mu_\theta(o_k),
 \operatorname{diag}(\sigma_\theta^2)\right) \nonumber\\
 &-\sum_{j=1}^{2}\log\!\left(1-\tanh^2(u_{k,j})\right) \nonumber\\
 &-\sum_{j=1}^{2}\log\!\left(\frac{a_j^{\max}-a_j^{\min}}{2}\right),
\end{align}
其中，$u_{k,j}$、$a_j^{\min}$ 和 $a_j^{\max}$ 分别是相应向量的第 $j$ 个分量。数值计算时，先用 $\varepsilon_{\mathrm{num}}>0$ 对每个 $1-\tanh^2(u_{k,j})$ 设置下界，再计算其对数。当前策略和采样策略采用相同的固定变换，因此这些 Jacobian 项会在重要性采样比中相消；上式给出计算有界动作对数概率和熵时使用的变换后概率密度。

完成一条 rollout 的数据采集后，广义优势估计（GAE）组合多个未来事件的时序差分残差，得到偏差与方差可调的策略优势估计 \cite{schulman2015gae}。令终止事件满足 $V_{\phi_{\mathrm{old}}}(o_K)=0$，则
\begin{align}
 \delta_k={}&R_k+\gamma V_{\phi_{\mathrm{old}}}(o_{k+1})
 -V_{\phi_{\mathrm{old}}}(o_k),\\
 \widehat{A}_k={}&\sum_{l=0}^{K-k-1}
 (\gamma\lambda_{\mathrm{GAE}})^l\delta_{k+l},\qquad
 \widehat{V}_k=\widehat{A}_k+V_{\phi_{\mathrm{old}}}(o_k).
\end{align}
其中，$\delta_k$ 是时序差分残差，$l$ 是非负的未来事件偏移量，$\lambda_{\mathrm{GAE}}$ 用于调节偏差与方差之间的权衡，$\widehat{A}_k$ 是优势估计，$\widehat{V}_k$ 是 critic 的训练目标。

基于变换后动作概率定义重要性采样比
\begin{equation}
 \rho_k(\theta)=
 \frac{\pi_\theta(a_k\mid o_k)}
 {\pi_{\theta_{\mathrm{old}}}(a_k\mid o_k)}.
\end{equation}
Actor 最大化带熵奖励的裁剪代理目标：
\begin{align}
 J_{\mathrm{actor}}(\theta)=\mathbb{E}_k\big[&
 \min\{\rho_k(\theta)\widehat{A}_k,
 \operatorname{clip}(\rho_k(\theta),1-\epsilon,1+\epsilon)
 \widehat{A}_k\}\nonumber\\
 &+\beta_H\mathcal{H}(\pi_\theta(\cdot\mid o_k))\big],
\end{align}
其中，$\mathbb{E}_k$ 表示对 rollout 中所存事件样本取经验平均；$\operatorname{clip}$ 将第一个参数截断到后两个参数指定的区间内；$\epsilon$ 是裁剪阈值；$\beta_H$ 是熵系数。变换后有界动作策略的微分熵定义为 $\mathcal{H}(\pi_\theta(\cdot\mid o_k))=-\mathbb{E}_{a\sim\pi_\theta(\cdot\mid o_k)}[\log\pi_\theta(a\mid o_k)]$。熵系数在训练过程中线性衰减，使随机探索在训练前期获得较高权重，并逐步减弱这一探索激励。critic 通过最小化下式进行训练：
\begin{equation}
 L_V(\phi)=\frac{1}{2}\mathbb{E}_k
 \left[\left(V_\phi(o_k)-\widehat{V}_k\right)^2\right].
\end{equation}

每条 rollout 保存观测、潜在动作和有界动作、奖励、采样策略对数概率以及价值估计。PPO 将这些事件样本划分为随机打乱的小批次，并在多个更新轮次中重复利用，对 $J_{\mathrm{actor}}$ 执行梯度上升、对 $L_V$ 执行梯度下降，随后再用更新后的策略采集新 rollout。确定性评估时，令 $u_k=\mu_\theta(o_k)$，并采用与训练阶段相同的双曲正切和仿射变换。

\section{案例研究与实验设置}
\subsection{数据来源与假设}
案例数据分为三类：公开事实、推导值和明确假设。公开事实包括线路名称、线路长度、车站数、车队规模、车辆容量，以及能够查到的储能技术。当没有准确的车站里程数据时，模型假设车站间距均匀。区间运行时间根据线路长度和公开的平均速度推导，超级电容能量则根据假设的电容和电压计算。其他假设包括各车站的乘客到达率曲线、上车服务时间、部分储能限制、辅助功率和共享总充电功率预算。本文不使用以往论文的数值结果来校准模型或设定验证阈值。

\thz{} 线路长 7.7 km，共 11 个车站、7 辆车，推导得到的超级电容可用容量为 4.734 kWh。\thp{} 线路长 14.3 km，共 19 个车站，高峰期使用 11 辆车；超级电容容量约为 4.66 kWh，并假设钛酸锂电池可用容量为 8.0 kWh。仿真时域为 07:30--09:30，即 $t_f-t_0=2$ h；初始时刻，各车辆被放置在各自运行循环的不同位置。模型假设两条线路的每个车站均可充电，并设置 $d^{\min}=10$ s、$d^{\max}=50$ s、$P^{\max}=300$ kW 和 $P^{\mathrm{off}}=5$ kW。

\subsection{实验因素}
实验在四个维度上展开：
\begin{itemize}
    \item 决策方法：w/o-Opt.、GA、PSO、DE、MPC 和 PPO；
    \item 扰动环境：无扰动以及四种稳健性环境 $E_1$、$E_2$、$E_3$ 和 $E_4$；
    \item 预算值：$B \in \{1600,1400,1200\}$ kW；
    \item 预算方式：两条线路共享总预算，或预先将总预算分给各线路。
\end{itemize}
无扰动环境令 $p_{\mathrm d}=0$。四种稳健性环境的参数对 $(p_{\mathrm d},\mu_{\mathrm d})$ 分别为 $E_1=(0.2,1.10)$、$E_2=(0.5,1.10)$、$E_3=(0.2,1.20)$ 和 $E_4=(0.5,1.20)$。对每个评估种子，两条线路分别使用由该种子确定的独立伪随机数流生成延误指示序列；在不同决策方法之间复用相同种子，以减少随机样本差异对方法比较的影响。分线路预算的主要对比设置为 $B_{\thz}=B_{\thp}=B/2$。

每一种 $B$ 与预算方式的组合都单独训练或优化。PPO、GA、PSO、DE 和 MPC 都使用修正后的实际时间功率超限积分和乘客等待时间目标。使用旧目标函数得到的结果不会在本次对比中重复使用。

\subsection{对比方法}
为了比较学习策略和显式短期前瞻，模型预测控制（MPC）基线反复优化由两条线路合并到站事件构成的有限预测时域，每次只执行第一个动作，并在下一事件重新求解。其内部预测器使用计划运行时间、两条线路的未来事件表、简化能耗模型，并假设当前总功率会作为背景功率持续存在。短期预测目标包含充电能量、功率超限和未满足能量项，并对未来乘客数进行近似。MPC 选出的每个动作最终都由完整仿真器按照与其他决策方法相同的事件成本进行评估。内部模型和未来事件信息使 MPC 相对于 PPO 具有明确的信息优势。

固定运行基线记为无优化（w/o-Opt.）。它在每个事件中采用相同的停站时间和请求充电功率，随后由仿真器根据可行性限制得到 $P_k$。为了比较在线反馈和整回合离线搜索，遗传算法（GA）、粒子群优化（PSO）和差分进化（DE）在相同动作范围内优化向量 $x=[d_1,\ldots,d_K,\widehat{P}_1,\ldots,\widehat{P}_K]$。向量中的每个元素都与线路代码、车辆编号和到访序号 $(\ell,i,v)$ 固定对应。因此，在不同扰动实验中，同一个向量元素（即基因）始终控制同一辆车的同一次到站。GA 使用交叉和变异操作，PSO 使用认知、社会和惯性项，DE 使用差分缩放和重组操作 \cite{holland1975adaptation,kennedy1995pso,storn1997de}。这些方法的具体参数与其他实验配置统一在本节报告。

固定策略和搜索类基线以同一成本定义最小化无折扣回合总成本 $\sum_{k=0}^{K-1}C_k$，从而保证各方法使用一致的评价口径。

\subsection{算法配置与评估方式}
表~\ref{tab:algorithm_settings} 汇总生成实验结果所采用的数值配置。PPO 进行确定性评估时，将潜在高斯分布的均值经过与训练阶段相同的双曲正切变换和仿射映射，并将观测归一化器固定为训练阶段得到的统计量。

\begin{table}[t]
\centering
\caption{实验采用的算法配置。}
\label{tab:algorithm_settings}
\footnotesize
\setlength{\tabcolsep}{3pt}
\begin{tabular}{lll}
\toprule
方法 & 参数 & 数值 \\
\midrule
PPO & 网络结构 & $3\times512$，Tanh \\
PPO & $\gamma/\lambda_{\mathrm{GAE}}/\epsilon$ & $0.993/0.98/0.2$ \\
PPO & 优化器/学习率 & Adam/$10^{-7}$ \\
PPO & 回合数/更新轮次/小批次 & $2000/64/64$ \\
PPO & 熵系数 & $10^{-2}\rightarrow0$（线性） \\
MPC & 预测时域 & 8 个合并到站事件 \\
w/o-Opt. & 停站时间/请求功率 & 30 s/300 kW \\
GA & 交叉率/变异率 & $0.9/0.1$ \\
PSO & 认知/社会/惯性系数 & $0.5/0.5/0.8$ \\
DE & 缩放因子/重组率 & $0.5/0.9$ \\
\bottomrule
\end{tabular}
\end{table}

\subsection{奖励权重设置}
所有参考对比实验均采用表~\ref{tab:cost_weights} 所列的固定成本权重。不同决策方法、扰动环境、功率预算、预算方式和随机种子使用完全相同的权重。这些设计系数预先继承自参考方法配置，并采用任意目标函数单位表示；其数值独立于试运行尺度和评估结果。

\begin{table}[t]
\centering
\caption{实验采用的参考成本权重。其中，cost 表示优化目标的任意数值单位。}
\label{tab:cost_weights}
\footnotesize
\setlength{\tabcolsep}{3pt}
\begin{tabular}{lccc}
\toprule
成本项 & 权重 & 数值 & 单位 \\
\midrule
实际存入能量 & $c_e$ & 5 & cost/kWh \\
功率预算超限 & $c_o$ & 10 & cost/(kW$^2$h) \\
未满足牵引能量 & $c_m$ & 1000 & cost/kWh \\
过充超额量 & $c_p$ & 750 & cost/kWh \\
乘客等待时间 & $c_w$ & 20 & cost/(乘客-分钟) \\
\bottomrule
\end{tabular}
\end{table}

\subsection{系数敏感性设计}
令 $c_r^{\mathrm{ref}}$ 表示表~\ref{tab:cost_weights} 中的参考权重。单因素敏感性实验对 $r\in\{o,w,m,p\}$ 取 $c_r=\beta_r c_r^{\mathrm{ref}}$，其中 $\beta_r\in\{0.5,1,2\}$，并保持其他权重为参考值。另一组权衡实验改变无量纲倍数之比 $\beta_o/\beta_w$；由于 $c_o$ 和 $c_w$ 对应的物理量不同，本文不把有量纲的 $c_o/c_w$ 解释为物理比值。每个实验点都在相同的随机种子和预算下重新训练或优化各决策方法。结果包括下文定义的乘客服务、未满足能量和功率超限指标。

\subsection{指标与图表}
每个训练或优化后的方案都使用随机种子 $\{0,1,2,3,4\}$ 进行评估。乘客服务指标包括：总等待时间（乘客-分钟）、每位已服务乘客的平均等待时间、仿真结束时仍在站台候车的人数、因车辆容量或停站时间限制而无法上车的人数，以及车辆载客率；载客率定义为车内乘客数除以车辆容量。功率预算指标包括：精确的超额功率平方积分（kW$^2$h）、功率超限持续时间（相关功率高于预算的总分钟数）和总功率峰值（kW）。采用分线路预算时，还会分别报告每条线路的超限持续时间和峰值功率。未满足牵引能量和各项成本也单独报告。每秒采样一次的功率曲线只用于绘图；所有功率超限指标都根据准确的充电开始和结束时刻计算。

决策方法对比实验报告以下内容：
\begin{itemize}
    \item $E_1$--$E_4$ 和预算配置下的 PPO 学习曲线；
    \item 带有预算参考线的 $P_{\thz}$、$P_{\thp}$ 和 $P_{\mathrm{tot}}$ 功率曲线；
    \item 比较共享预算和分线路预算的权衡曲线；
    \item 所有决策方法的成本成分表；
    \item 跨线路协调实例，例如 \thz{} 处于充电高峰时，\thp{} 选择暂不充电。
\end{itemize}

\section{结果与讨论}
本文研究运营商应如何在多条有轨电车线路之间分配共同的合同需量上限或已购买的功率容量，以及这种分配如何影响乘客服务和车辆满足牵引能量需求的能力。因此，共享预算实验结果应理解为运营规划层面的对比。

另一个局限是，乘客到达率曲线和若干储能参数属于明确列出的建模假设。因此，这些实验代表本文定义的工程场景，所得结论适用于两条广州有轨电车线路的模拟运行。模型可信度主要来自前后一致的假设、固定随机种子下可复现的结果，以及不同扰动和预算条件下的敏感性分析。

\section{结论}
本文首次建立事件驱动模型，联合控制多条无接触网有轨电车线路的停站时间和充电功率。该框架在一个离散事件仿真与强化学习模型中，同时考虑不同时间发生的到站事件、不同的车载储能技术、乘客服务和运营商共享的总充电功率预算。广州两线路案例表明，在线 DRL 决策智能体可以在运行时间受到扰动时减少同时充电造成的功率峰值和预算超限，同时保证车辆有足够能量运行并维持乘客服务质量。其综合表现优于固定运行、种群优化和滚动时域控制。这些结果表明，自适应反馈控制有助于协调多条有轨电车线路的功率需求和实时运行。

\bibliographystyle{IEEEtran}
\bibliography{references}

\end{CJK*}
\end{document}
